% The pre-chapters
\chapter*{Abstract} %pre-chapters should not be numbered, hence the "*"
\pagenumbering{roman} %introductory pages should be roman
\setcounter{page}{1}

\addcontentsline{toc}{chapter}{\protect\numberline{}Abstract} %add the chapter to the table of contents, this is not automatically added when creating unnumbered chapters (*). Add it in a chapter style, and keep all chapters on the same numberline indent regardless of number or not on the chapter
Geospatial datasets have grown large enough that downloading entire files to answer a query is often impractical. Cloud-optimized formats address this by keeping data on object storage and reading only the byte ranges a task needs. On pay-as-you-go cloud infrastructure, where network transfer is billed by the byte and compute by the hour, performance and operational cost no longer move together. This thesis benchmarks cloud-native and traditional geospatial technologies on common Microsoft Azure infrastructure, using identical Norwegian building-polygon data. It measures wall-clock run-time, network transfer, and operational cost across representative vector workloads. Four configurations are compared: DuckDB over GeoParquet, PostGIS, GeoPandas over Shapefile, and Apache Sedona on Databricks. The single-machine workloads cover point-in-polygon, k-nearest-neighbor, and bounding-box queries, while a national-scale spatial join sets the distributed engine against the single-node engines across small, medium, and large dataset tiers.

On a single machine, no configuration was superior on every axis. A warm, indexed PostGIS held the lowest run-time in every single-machine cell. The decoupled read model of GeoParquet and DuckDB yielded no run-time advantage at the sub-second scale of these queries. Network transfer was a per-architecture signature rather than a single ranking, because the bytes crossing the host boundary denote a scanned input for DuckDB and a returned result set for PostGIS. At the scale tested, this transfer did not contribute to operational cost, since the egress term was zero for every single-node configuration. Even so, DuckDB recorded the lowest cost on two of the small-tier query patterns.

For the national-scale spatial join, the answer hinged on the size of the join. At the large tier, neither single-node engine completed the join while the distributed broadcast join on Apache Sedona did, making distribution a matter of feasibility rather than speed. At the medium tier the broadcast join was faster than the fastest single-node engine by close to an order of magnitude or more at every worker count. This difference is carried by a complete effect size that the fixed three-iteration budget nonetheless leaves formally non-significant. For a small join a tuned single node was not overtaken across the 2--16 workers tested. The choice of join strategy was the largest single lever, with the partitioned join running roughly two orders of magnitude slower than the broadcast join where it completed. Read together, these orderings held no consistent winner: the leading configuration depended jointly on the query pattern, the dataset tier, and the outcome dimension.

The work contributes the first reported comparison of a distributed engine against single-node PostGIS and DuckDB on a national-scale spatial join. It also reports network transfer jointly with run-time and operational cost, alongside effect sizes and bootstrapped confidence intervals. The orderings that follow from format and engine structure are expected to reach beyond the tested conditions, while the absolute times and costs, priced from provider- and region-specific rates, are not.

\chapter*{Sammendrag}

Geospatiale datasett har blitt så store at det ofte er upraktisk å laste ned hele filer for å besvare en spørring. Skyoptimaliserte formater løser dette ved å la dataene ligge i objektlagring og bare lese de byteområdene en oppgave trenger. På skyinfrastruktur med forbruksbasert betaling, der nettverksoverføring faktureres per byte og beregningskraft per time, varierer ytelse og driftskostnad ikke lenger i takt. Denne oppgaven benchmarker cloud-native og tradisjonelle geospatiale teknologier på felles Microsoft Azure-infrastruktur med identiske norske bygningspolygondata, og måler kjøretid, nettverksoverføring og driftskostnad på tvers av representative vektorarbeidslaster. Fire konfigurasjoner sammenlignes: DuckDB over GeoParquet, PostGIS, GeoPandas over Shapefile og Apache Sedona på Databricks. Arbeidslastene på én maskin omfatter punkt-i-polygon-, k-nærmeste-nabo- og bounding-box-spørringer, mens en spatial join i nasjonal skala stiller den distribuerte motoren mot enkeltnode-motorene over datasettnivåene liten, middels og stor.

På én enkelt maskin var ingen konfigurasjon best på alle akser. En varm, indeksert PostGIS hadde lavest kjøretid i hver celle for enkeltmaskinarbeidslastene. Den frakoblede lesemodellen til GeoParquet og DuckDB ga ingen kjøretidsfordel på subsekundskalaen til disse spørringene. Nettverksoverføring var en arkitekturspesifikk signatur snarere enn én rangering: bytene som krysser vertsgrensen, utgjør skannet inndata for DuckDB og et returnert resultatsett for PostGIS. På den testede skalaen bidro overføringen ikke til driftskostnaden, ettersom egress-leddet var null for hver enkeltnode-konfigurasjon. Likevel hadde DuckDB lavest kostnad for to av spørringsmønstrene på det minste nivået.

For spatial join i nasjonal skala var svaret avhengig av størrelsen på joinen. På det største nivået fullførte ingen av enkeltnode-motorene joinen, mens den distribuerte broadcast join-en på Apache Sedona gjorde det — distribusjon ble dermed et spørsmål om gjennomførbarhet snarere enn hastighet. På det middels nivået var broadcast join-en raskere enn den raskeste enkeltnode-motoren med nær en størrelsesorden eller mer ved hvert antall workere, en forskjell understøttet av en fullstendig effektstørrelse som det faste budsjettet på tre iterasjoner likevel gjør formelt ikke-signifikant. For en liten join ble en finjustert enkeltnode ikke forbigått over de 2--16 testede workerne. Valget av join-strategi var den enkeltfaktoren som betydde mest, ettersom den partisjonerte joinen kjørte omtrent to størrelsesordener saktere enn broadcast join-en der den fullførte. Lest sammen ga disse rangeringene ingen konsistent vinner: den ledende konfigurasjonen avhang av kombinasjonen av spørringsmønster, datasettnivå og utfallsdimensjon.

Arbeidet bidrar med den første rapporterte sammenligningen av en distribuert motor opp mot enkeltnode-PostGIS og DuckDB på en spatial join i nasjonal skala. Det rapporterer også nettverksoverføring sammen med kjøretid og driftskostnad, i tillegg til effektstørrelser og bootstrappede konfidensintervaller. Rangeringene som følger av format- og motorstruktur, forventes å gjelde utover de testede betingelsene, mens de absolutte tidene og kostnadene, beregnet ut fra leverandør- og regionspesifikke priser, ikke gjør det. %insert the chapter text from the files

\chapter*{Preface}
\addcontentsline{toc}{chapter}{\protect\numberline{}Preface} 
\begin{quote}
    \itshape\large
    ``Talk is cheap. Show me the code.''
    \begin{flushright}
        \normalfont\normalsize --- Linus Torvalds, 2000
    \end{flushright}
\end{quote}


\cleardoublepage
\phantomsection
\addcontentsline{toc}{chapter}{\protect\numberline{}\contentsname}
\tableofcontents

\cleardoublepage
\phantomsection
\addcontentsline{toc}{chapter}{\protect\numberline{}\listfigurename}
\listoffigures

\cleardoublepage
\phantomsection
\addcontentsline{toc}{chapter}{\protect\numberline{}\listtablename}
\listoftables

\cleardoublepage

\chapter*{Abbreviations}
\addcontentsline{toc}{chapter}{\protect\numberline{}Abbreviations}

% The acronym list is a page-spanning longtable. A single \markboth before it
% sets the running mark on modern TeX, but the mark is not reliably re-emitted
% on the table's continuation pages across engines, so those pages can inherit
% the previous section's mark and print "List of Tables". Drive the header from
% a dedicated page style with a literal title instead of \leftmark, so every
% continuation page reads "Abbreviations" regardless of mark propagation. The
% \markboth is kept for the first page's mark and PDF tooling.
\markboth{\MakeUppercase{Abbreviations}}{\MakeUppercase{Abbreviations}}
\fancypagestyle{abbrevpage}{%
	\fancyhf{}%
	\renewcommand{\headrulewidth}{0pt}%
	\fancyhead[LE,RO]{\thepage}%
	\fancyhead[RE,LO]{\MakeUppercase{Abbreviations}}%
}
\pagestyle{abbrevpage}

\begingroup
\renewcommand*{\glossarysection}[2][]{}
\printglossary[type=\acronymtype]
\endgroup

\cleardoublepage